\begin{table}[!htbp]
\centering
\small
\caption{SEER-HCC income analysis: adjusted unstratified Cox model}
\begin{tabular}{llll}
\toprule
Variable & HR & 95\% CI & P-value \\
\midrule
Income Group (High) & 0.750 & (0.677, 0.832) & \textless{}0.001 \\
Age & 1.013 & (1.008, 1.018) & \textless{}0.001 \\
Race (Black) & 0.898 & (0.774, 1.041) & 0.154 \\
Race (Asian or Pacific Islander) & 0.847 & (0.728, 0.986) & 0.032 \\
Sex (Male) & 1.099 & (0.977, 1.236) & 0.114 \\
T Stage (T2) & 1.356 & (1.175, 1.564) & \textless{}0.001 \\
T Stage (T3) & 2.594 & (2.264, 2.972) & \textless{}0.001 \\
T Stage (T4) & 2.856 & (2.452, 3.326) & \textless{}0.001 \\
N Stage (N1) & 1.382 & (1.165, 1.640) & \textless{}0.001 \\
M Stage (M1) & 2.036 & (1.743, 2.378) & \textless{}0.001 \\
Tumor Size & 1.004 & (1.003, 1.004) & \textless{}0.001 \\
Tumor Grade (II) & 0.991 & (0.875, 1.122) & 0.886 \\
Tumor Grade (III) & 1.117 & (0.960, 1.299) & 0.151 \\
AFP Status (Positive) & 1.674 & (1.496, 1.872) & \textless{}0.001 \\
\bottomrule
\end{tabular}
\begin{minipage}{0.95\linewidth}
\vspace{0.5em}
\footnotesize N = 4668; events = 1608. HR: adjusted hazard ratio; CI: confidence interval. P-values are Wald tests. Categorical reference levels: Income Group = Low; Race = White; Sex = Female; T Stage = T1; N Stage = N0; M Stage = M0; Tumor Grade = I; AFP Status = Negative. Age and tumor size HRs are per year and per mm, respectively.
\end{minipage}
\end{table}
